\documentclass[11pt,a4paper,UTF8,fontset=none]{ctexart}
\setCJKmainfont[
  Path=assets/fonts/,
  AutoFakeBold=2.2
]{NotoSansSC.ttf}
\setCJKsansfont[
  Path=assets/fonts/,
  AutoFakeBold=2.2
]{NotoSansSC.ttf}
\setCJKmonofont[
  Path=assets/fonts/
]{NotoSansSC.ttf}
\usepackage[a4paper,margin=2.25cm]{geometry}
\usepackage{booktabs,longtable,array,tabularx}
\usepackage{graphicx}
\usepackage{xcolor}
\usepackage{amsmath,amssymb}
\usepackage{listings}
\usepackage{hyperref}
\usepackage{fancyhdr}
\usepackage{microtype}
\usepackage{enumitem}
\usepackage{caption}
\usepackage{tikz}
\usepackage{float}
\usepackage{placeins}
\usetikzlibrary{arrows.meta,positioning,shapes.geometric,calc}

\definecolor{navy}{HTML}{16324F}
\definecolor{blue}{HTML}{2D6A9F}
\definecolor{lightblue}{HTML}{EAF2F8}
\definecolor{warning}{HTML}{9C2F2F}
\definecolor{flowgreen}{HTML}{E8F3EC}
\definecolor{flowyellow}{HTML}{FFF4D6}
\hypersetup{colorlinks=true,linkcolor=navy,urlcolor=blue,citecolor=navy}
\setlist{nosep,leftmargin=2em}
\setlength{\parindent}{2em}
\setlength{\parskip}{0.25em}
\setlength{\emergencystretch}{2em}
\renewcommand{\arraystretch}{1.15}
\pagestyle{fancy}
\fancyhf{}
\fancyhead[L]{对称密码的软件实现与优化}
\fancyhead[R]{课程项目}
\fancyfoot[C]{\thepage}
\setlength{\headheight}{14pt}
\lstset{
  basicstyle=\ttfamily\small,
  keywordstyle=\color{blue},
  commentstyle=\color{gray},
  frame=single,
  breaklines=true,
  columns=fullflexible,
  backgroundcolor=\color{lightblue},
  xleftmargin=1em,
  xrightmargin=1em
}
\tikzset{
  flowproc/.style={
    draw=navy, fill=lightblue, rounded corners=2pt, align=center,
    minimum height=9mm, text width=31mm, inner sep=4pt, font=\small
  },
  flowio/.style={
    draw=blue, fill=flowgreen, rounded corners=8pt, align=center,
    minimum height=9mm, text width=28mm, inner sep=4pt, font=\small
  },
  flowdecision/.style={
    diamond, aspect=2.2, draw=warning, fill=flowyellow, align=center,
    inner sep=1.5pt, font=\small
  },
  flowerror/.style={
    draw=warning, fill=red!7, rounded corners=2pt, align=center,
    minimum height=8mm, text width=27mm, inner sep=3pt, font=\small
  },
  flowarrow/.style={-{Stealth[length=2.2mm]}, thick, draw=navy},
  flowaux/.style={-{Stealth[length=2mm]}, thick, dashed, draw=blue}
}

\input{tmp/generated_results.tex}

\title{\bfseries 对称密码算法的软件实现与优化\\[0.4em]
\large AES／SM4／GIFT／TWINE 与 CTR／GCM／XTS}
\author{课程项目实现与实验报告}
\date{2026 年 7 月}

\begin{document}
\maketitle

\begin{abstract}
本项目从统一的 C11 标量基线出发，实现 AES-128/192/256、SM4-128、
GIFT-64（128-bit key）和 TWINE-80/128，并系统比较 T-table、
shuffle/bitslice/fixslice 与新指令集。工作模式覆盖全部算法的 CTR，以及
AES/SM4 的 GCM、XTS；GCM 使用 ARM PMULL 加速 GHASH，XTS 同时支持
IEEE~1619 与 GB/T~17964-2021 的 tweak 乘法约定及 ciphertext stealing。
Apple M2 Pro 与 Intel Core i9-13900H 上分别保存 3 次预热、15 次正式测量
的全部样本。x86-64 原生验证执行 AES-NI、VAES、SSSE3、GFNI、PCLMUL 和
VPCLMUL，包含 \XGenericChecks 项通用检查、\XBackendChecks 项专用后端检查
与 \XOpenSSLChecks 项 OpenSSL 差分。Intel VSM4 后端已完成适配和静态指令
检查，但实验使用的 Core i9-13900H 不支持该指令集，因此没有 VSM4 原生
性能数据。
\end{abstract}

\noindent\textbf{关键词：} 对称密码；AES；SM4；GIFT；TWINE；T-table；
NEON TBL；PMULL；GFNI；CTR；GCM；XTS

\clearpage
\tableofcontents
\clearpage

\section{实现顺序与交付边界}

项目按照“先建立正确性基线，再逐层替换热点”的顺序实现：
\begin{enumerate}
  \item 先完成 AES、SM4、GIFT 与 TWINE 的单块加/解密和密钥扩展；
  \item 用官方或论文向量固定参考输出，形成所有优化后端的比较基线；
  \item 为 AES/SM4 加入 T-table，对比表尺寸、缓存占用与吞吐率；
  \item 加入 NEON TBL、x86 PSHUFB 和 GIFT bitslice，并把单块接口扩展为
  四路/八路批量内核；
  \item 接入 AES、GFNI、PMULL/PCLMUL 等指令，通过运行时特征检测选择
  可执行后端；
  \item 最后让 CTR、GCM 与 XTS 调用批量分组接口，并补充边界测试、
  OpenSSL 差分、sanitizer 和六档长度基准。
\end{enumerate}

后续章节按同一顺序展开具体实现：先说明如何把 AES/SM4 的非线性层与线性层
合并成 T-table，再说明 shuffle/bitslice 与密码指令怎样替换查表；模式部分
则直接结合 \texttt{src/modes.c}，给出 CTR 的计数器批量生成、GCM 的
CTR+GHASH 组合以及 XTS 的 tweak/XEX/CTS 流程。所有代码片段都由项目中的
实际实现精简而来，而不是独立的演示伪代码。

整个实现始终保持三条不变量。第一，参考后端定义字节序、轮密钥次序和错误
语义，优化后端只能改变内部数据表示，不能改变外部密文。第二，单块接口和
多块接口共享同一 \texttt{sc\_ctx}，批量宽度不足时只回退剩余分组，而不会
重新扩展密钥。第三，工作模式只依赖统一分组接口，不直接识别 AESE、PSHUFB
或 GFNI；因此替换后端后，CTR/GCM/XTS 可以自动获得批量加密能力，同时仍由
模式层负责计数器、认证和 ciphertext stealing 的正确性。

工程上把“功能完成”和“特定设备可执行”分开记录。源码可以为某条目标 ISA
生成正确指令，但只有运行时特征检测通过后才会进入该后端。这样既能保留
ARM64 与 x86-64 的同一套公共 API，也不会把静态编译成功误写成原生性能
结果。本文所有性能结论均对应实际保存的 ARM 或 x86 测量样本。

\begin{quote}\color{warning}
\textbf{安全边界：}本项目是可复现实验实现，不是经过认证的生产密码库。尤其
T-table 会以秘密值索引缓存，不应成为真实系统的默认后端。
\end{quote}

\section{微架构基础与优化模型}

\subsection{吞吐率不是指令条数}

现代处理器可多发射并行执行，实际耗时由执行端口、依赖链、指令时延和
吞吐共同决定。因而同样的轮函数可以有三类改进：
\begin{enumerate}
  \item 缩短关键依赖链，例如把互不相关的分组交错执行；
  \item 减少内存层次开销，例如把状态保持在寄存器，或缩小表；
  \item 用一条 SIMD/密码指令替代多条标量指令。
\end{enumerate}
全展开可以删除循环控制，却会扩大指令缓存占用。本项目保留清晰的循环基线，
只在数据表示和后端处做可测量的变换，同时以对象尺寸 CSV 记录代码代价。

分组密码的单个分组往往存在“第 \(i+1\) 轮依赖第 \(i\) 轮”的长链，单纯
增加 SIMD 指令并不能消除这条依赖。本项目因此把互不相关的四个或八个分组
放入同一批次：当一个分组等待前一条指令结果时，处理器仍可发射其他分组的
轮指令。AES 硬件后端、SM4 shuffle/GFNI 和工作模式的计数器生成都采用这一
思路。批量宽度是实现选择，不是 API 限制；调用者给出任意分组数，后端先
处理完整批次，再用较窄内核或标量路径完成余数。

缓存行为同样影响结果。4~KiB T-table 可能全部驻留 L1，但仍会与轮密钥、
输入输出和其他代码竞争；1~KiB 表减少数据占用，却增加旋转指令；固定扫描
shuffle 不按秘密值选择地址，但每轮需要读取或组合更多表行。报告因此同时
记录吞吐、表大小和对象大小，不用“指令更少”替代实际测量。

\subsection{端序与未定义行为}

所有外部字节串都通过 \texttt{sc\_load\_be32/64} 与
\texttt{sc\_store\_be32/64} 显式转换。2~KiB SM4 重叠表虽然利用偏移
读取旋转后的 32-bit 字，却不把非对齐地址强制转换成 \texttt{uint32\_t*}，
从而避免别名和对齐未定义行为。这一点由 ASan/UBSan 与
\texttt{-Werror} 共同约束。

AES 和 SM4 的规范都以字节串描述输入，但二者内部字排列不同；GIFT/TWINE
又以 bit 或 nibble 为基本单位。实现只在加载和保存边界进行显式转换，中间
轮函数始终操作固定宽度无符号整数。移位量均小于类型宽度，旋转由左右移和
按位或构成；所有可能非对齐的输入先复制到局部数组。因此原地操作、
非对齐地址和不同编译器优化级别不会依赖宿主机未规定的行为。

\section{四种分组密码的基础实现}

\subsection{AES}

AES 按 FIPS~197 实现 128/192/256-bit 密钥扩展、SubBytes、ShiftRows、
MixColumns 及逆变换。参考路径以 16 字节状态表示；T-table 同时实现正向与
逆向四表。ARM 路径用 AESE/AESMC 与 AESD/AESIMC，x86 路径用 AES-NI，
多块时可进入 VAES。所有路径使用同一轮密钥并与 FIPS KAT 交叉比较。

输入按列映射为 \(4\times4\) 字节状态。加密先异或第 0 轮密钥，中间轮依次
执行 SubBytes、ShiftRows、MixColumns 和 AddRoundKey，末轮省略
MixColumns；AES-128/192/256 分别执行 10、12、14 轮。密钥扩展把原始密钥
扩展为 \(N_r+1\) 组 128-bit 轮密钥，并根据 \(N_k=4,6,8\) 分别处理
RotWord、SubWord、Rcon 和 AES-256 的额外 SubWord 分支。

解密参考路径按相反顺序执行逆变换。为了让 AESD/AESIMC、AES-NI AESDEC
以及逆 T-table 都直接消费合适的轮密钥，初始化阶段同时准备逆序轮密钥，
并对中间轮密钥执行 InvMixColumns。多块接口加载若干独立状态到寄存器，
每轮对全部状态应用同一轮密钥；末尾不足并行宽度的分组交给窄批次处理。
这样密钥扩展不进入稳态计时，也避免“硬件名称后端实际逐块调用参考代码”。

\subsection{SM4}

SM4 使用 32 轮广义 Feistel 更新：
\[
  X_{i+4}=X_i\oplus L\!\left(\tau(X_{i+1}\oplus X_{i+2}
  \oplus X_{i+3}\oplus rk_i)\right),
\]
\[
  L(B)=B\oplus(B\lll2)\oplus(B\lll10)\oplus(B\lll18)\oplus(B\lll24).
\]
密钥扩展和数据轮均依据 GB/T~32907-2016。实现把输入、输出与轮密钥端序
明确分开，并用官方向量
\texttt{012345...3210 $\rightarrow$ 681edf...4246} 验证。

数据路径只需维护四个 32-bit 滚动状态字。每轮先计算
\(T=X_{i+1}\oplus X_{i+2}\oplus X_{i+3}\oplus rk_i\)，再执行逐字节
S 盒 \(\tau\) 和线性层 \(L\)，最后与 \(X_i\) 异或。第 32 轮完成后以
\(X_{35},X_{34},X_{33},X_{32}\) 的逆序写出。滚动窗口减少了长期存活的
临时量，批量后端则为每个 lane 保留独立四字状态，使多块之间没有数据依赖。

密钥扩展先将 128-bit 主密钥与系统参数 \(FK_0\ldots FK_3\) 异或，再对
32 个 \(CK_i\) 常量执行与数据轮类似的非线性变换，但使用
\[
 L'(B)=B\oplus(B\lll13)\oplus(B\lll23).
\]
产生的 32 个轮密钥按正序用于加密，按逆序用于解密，因此 SM4 的解密无需
另写一套轮函数。T-table、shuffle、GFNI、AES 辅助和专用指令适配都遵守同一
轮密钥次序，并在测试中同时比较加密和解密输出。

\subsection{GIFT-64 与 TWINE}

GIFT-64 采用四路 bitslice 状态，S 盒直接变为四个位平面上的布尔网络，
位置换由固定掩码网络完成；批量 API 每次处理四个独立分组。TWINE 以
16 个 nibble 表示状态，分别实现 80-bit 与 128-bit
密钥调度；KAT 密文为 \texttt{7c1f0f80b1df9c28} 与
\texttt{979ff9b379b5a9b8}。

GIFT-64 的 64-bit 状态在 28 轮中依次经过 4-bit S 盒、固定 bit
permutation 和轮密钥/轮常量异或。独立参考实现按 nibble 处理一个分组，
用于确认轮常量、bit 编号和密钥更新；优化实现把四个分组的同一 bit 位置
转置到四个位平面，再以 AND、XOR、NOT 表达 S 盒。固定置换被编译为掩码、
移位和交换网络，解密使用逆布尔网络与逆置换，四个 lane 最后再转回字节序。

TWINE 的 64-bit 状态由 16 个 nibble 构成，共执行 36 轮。每轮只对偶数
nibble 执行“轮密钥异或、S 盒、异或到相邻 nibble”，随后执行固定置换；
最后一轮不再置换。80-bit 与 128-bit 版本的区别主要在密钥状态宽度和轮密钥
抽取/更新位置，数据轮可以共用。shuffle 后端把 16 个 nibble 打包进
128-bit 向量，第一次 TBL/PSHUFB 完成 S 盒查找，第二次重排完成轮置换，
从而让多块数据布局和算法的 nibble 结构一致。
\section{统一接口与运行时分派}

公共上下文 \texttt{sc\_ctx} 对外保持不透明，提供密钥初始化、单块/多块
加解密、块长/算法/后端查询和明确错误码。单块 API 支持原地操作。
后端可选 \texttt{ref}、三种 \texttt{ttable}、\texttt{shuffle}、
\texttt{aes-hw}、\texttt{gfni}、\texttt{sm4-hw} 与 \texttt{auto}。

\begin{lstlisting}[language=C,caption={统一接口的核心调用}]
sc_ctx ctx;
sc_status rc = sc_init(&ctx, SC_AES_128, SC_BACKEND_AUTO,
                       key, sizeof(key));
rc = sc_encrypt_blocks(&ctx, input, output, blocks);
\end{lstlisting}

ARM64 检测 AES、PMULL 与 SM4 能力；x86 检测 AES-NI、SSSE3、AVX2、
VAES、PCLMUL、VPCLMUL 和 GFNI。显式请求不可用 ISA 时返回
\texttt{SC\_ERR\_UNSUPPORTED}，不做静默冒充。

\texttt{sc\_init} 首先校验算法与密钥长度，再完成一次密钥扩展和后端选择。
\texttt{auto} 按“可用的硬件/常量访问后端优先、参考后端兜底”的规则选择；
显式后端则必须同时满足算法适用性和 CPU 特征。例如 AES 不能选择 GFNI，
GIFT/TWINE 不能选择 T-table，缺少 VSM4 的处理器不能初始化
\texttt{sm4-hw}。最终实际选择可由 \texttt{sc\_backend\_id} 查询，便于
测试和基准确认没有发生静默回退。

\texttt{sc\_encrypt\_blocks} 与 \texttt{sc\_decrypt\_blocks} 根据
上下文直接进入批量内核。AES 硬件路径、SM4 shuffle/GFNI/AES 辅助路径、
GIFT bitslice 和 TWINE shuffle 都在这里处理多块；只有尾端不足 lane 宽度
时才调用较窄路径。API 允许输入输出地址相同，并在进入后端前检查上下文、
指针和算法状态。CTR/GCM/XTS 因而只需按批次组织输入，不需要复制每种算法
的密钥调度或轮函数。

\section{优化技术}

\begin{figure}[H]
\centering
\begin{tikzpicture}
  \node[flowio, text width=34mm] (ref) at (0,0)
    {标量参考轮函数\\逐步执行 S 盒、线性层和置换};
  \node[flowproc, text width=45mm] (hot) at (0,-1.9)
    {分析热点与并行性\\确定查表、数据布局和 ISA 替换位置};

  \node[flowproc, text width=32mm] (tt) at (-6,-4.6)
    {\textbf{T-table}\\预计算 S 盒+线性层\\1/2/4~KiB 表查找并异或};
  \node[flowproc, text width=32mm] (shuf) at (-2,-4.6)
    {\textbf{Shuffle}\\nibble/固定 16 行表\\TBL/PSHUFB 多块并行};
  \node[flowproc, text width=32mm] (slice) at (2,-4.6)
    {\textbf{Bitslice/fixslice}\\转置多个分组\\布尔 S 盒与跨轮置换合并};
  \node[flowproc, text width=32mm] (isa) at (6,-4.6)
    {\textbf{硬件指令}\\运行时特征检测\\AES/GFNI/PMULL 批量内核};

  \node[flowproc, text width=48mm] (multi) at (0,-7.6)
    {统一多块后端\\
     \texttt{sc\_encrypt\_blocks}\\
     \texttt{sc\_decrypt\_blocks}};
  \node[flowio, text width=43mm] (modes) at (0,-10.0)
    {复用到 CTR/GCM/XTS\\处理计数器、GHASH、tweak 与尾部};

  \draw[flowarrow] (ref) -- (hot);

  % 上方分流总线：先水平展开，再分别垂直进入四个后端，避免斜线穿过节点。
  \coordinate (split) at (0,-3.15);
  \coordinate (splitL) at (-6,-3.15);
  \coordinate (splitR) at (6,-3.15);
  \draw[thick,draw=navy] (hot.south) -- (split);
  \draw[thick,draw=navy] (splitL) -- (splitR);
  \draw[flowarrow] (splitL) -- (tt.north);
  \draw[flowarrow] (-2,-3.15) -- (shuf.north);
  \draw[flowarrow] (2,-3.15) -- (slice.north);
  \draw[flowarrow] (splitR) -- (isa.north);

  % 下方汇流总线：各后端先独立落到总线，再由单一箭头进入批量 API。
  \coordinate (joinL) at (-6,-6.35);
  \coordinate (joinR) at (6,-6.35);
  \coordinate (join) at (0,-6.35);
  \draw[thick,draw=navy] (tt.south) -- (joinL);
  \draw[thick,draw=navy] (shuf.south) -- (-2,-6.35);
  \draw[thick,draw=navy] (slice.south) -- (2,-6.35);
  \draw[thick,draw=navy] (isa.south) -- (joinR);
  \draw[thick,draw=navy] (joinL) -- (joinR);
  \draw[flowarrow] (join) -- (multi.north);
  \draw[flowarrow] (multi) -- (modes);
\end{tikzpicture}
\caption{从标量轮函数到 T-table、shuffle、bitslice/fixslice 和硬件指令后端的优化流程。}
\end{figure}

\subsection{T-table：空间换指令}

T-table 的实现分两步。第一步在初始化时枚举全部 256 个字节，将 S 盒输出和
线性变换预先合并；第二步在每一轮中从状态的四个字节分别查表并异或轮密钥。
以 AES 为例，若 $s=S(i)$，则首表项包含
$\{02\}s,\{01\}s,\{01\}s,\{03\}s$，其余三张表由字节旋转得到。这样一次
查表异或同时完成 SubBytes、ShiftRows 对应的取字节和 MixColumns。

设输入状态按四个大端字 \(s_0,s_1,s_2,s_3\) 保存，则一个输出字由
\[
 T_0[s_0^{(0)}]\oplus T_1[s_1^{(1)}]\oplus
 T_2[s_2^{(2)}]\oplus T_3[s_3^{(3)}]\oplus rk
\]
得到。四个索引分别来自 ShiftRows 后落入同一列的字节，因此不需要先把
16 字节状态显式重排。每张表有 256 个 32-bit 项，共占 4096 字节；四个
输出字可以交错计算，以减少单一查表结果形成的依赖链。表仅保存与密钥无关
的 S 盒和线性变换，轮密钥仍在每轮末尾异或。

\begin{lstlisting}[language=C,caption={AES T-table 的构造与一列轮更新（src/aes.c）}]
for (unsigned i = 0; i < 256; ++i) {
    uint8_t s = aes_sbox[i];
    uint32_t x = ((uint32_t)gf_mul(s, 2) << 24) |
                 ((uint32_t)s << 16) | ((uint32_t)s << 8) |
                 gf_mul(s, 3);
    te[0][i] = x;
    te[1][i] = (x >> 8)  | (x << 24);
    te[2][i] = (x >> 16) | (x << 16);
    te[3][i] = (x >> 24) | (x << 8);
}

t[0] = te[0][s0 >> 24] ^ te[1][(s1 >> 16) & 255] ^
       te[2][(s2 >> 8) & 255] ^ te[3][s3 & 255] ^ rk[0];
\end{lstlisting}

末轮不能包含 MixColumns，因此仍单独执行 S 盒、ShiftRows 和轮密钥异或；
解密方向建立逆 S 盒与 InvMixColumns 合并的 \texttt{td[4][256]}，并对中间
轮密钥应用逆列混合。SM4 的表项则预计算
$L(S(i)\ll24)$。四表方案保存它的三个旋转版本；1~KiB 方案只保存首表并在
轮内旋转；2~KiB 方案把每个 32-bit 表项连续复制两次，从偏移 3、2、1
读取旋转后的四字节序列。

SM4 的轮输入 \(T\) 被拆成 \(b_0,b_1,b_2,b_3\) 四个字节。4~KiB 方案直接
计算 \(T_0[b_0]\oplus T_1[b_1]\oplus T_2[b_2]\oplus T_3[b_3]\)；
1~KiB 方案从同一表取四个值后旋转 0、8、16、24 bit；2~KiB 方案利用每项
连续复制的八个字节，让不同起始偏移读出的四字节窗口等价于旋转结果。三条
路径的输出必须逐字等于标量 \(L(\tau(T))\)，解密只需把轮密钥反向后复用
同一查表轮函数。

\begin{lstlisting}[language=C,caption={SM4 三种 T-table 数据路径（src/sm4.c）}]
uint32_t a = sm4_t[0][b0], b = sm4_t[0][b1];
uint32_t c = sm4_t[0][b2], d = sm4_t[0][b3];

/* 1 KiB: one table plus rotations */
return a ^ ((b >> 8) | (b << 24)) ^
       ((c >> 16) | (c << 16)) ^ ((d >> 24) | (d << 8));

/* 2 KiB: duplicated entries permit overlapping byte loads */
return load32(overlap[b0])     ^ load32(overlap[b1] + 3) ^
       load32(overlap[b2] + 2) ^ load32(overlap[b3] + 1);

/* 4 KiB: four independent 256-entry tables */
return sm4_t[0][b0] ^ sm4_t[1][b1] ^
       sm4_t[2][b2] ^ sm4_t[3][b3];
\end{lstlisting}

\begin{table}[htbp]
\centering
\caption{表布局与数据路径}
\begin{tabular}{lll}
\toprule
算法/后端 & 表大小 & 每轮方法\\
\midrule
AES 四表 & 4~KiB & SubBytes+ShiftRows+MixColumns 合并\\
SM4 四表 & 4~KiB & 四个字节分别查表后异或\\
SM4 单表 & 1~KiB & 查一个表，再旋转 8/16/24 bit\\
SM4 重叠 & 2~KiB & 8 字节重复项，偏移 0/3/2/1 读取\\
\bottomrule
\end{tabular}
\end{table}

表越小不保证越快：1~KiB 路径增加旋转，2~KiB 路径增加非对齐字节组装；
4~KiB 路径则给 L1 数据缓存更大压力。四条路径同时存在，性能结论交给同一
基准而不是先验判断。

T-table 的优化边界也很明确：它减少逻辑运算，却把由状态字节决定的索引
暴露给缓存层次。攻击者如果能观察 cache set、预取或执行时间，可能推断
密钥相关信息。因此运行时不会把“表路径最快”自动等价为“最安全”；报告把
它作为一种数据路径单独测量，并与固定访问的 shuffle、bitslice 和硬件密码
指令比较。

\subsection{Shuffle、bitslice 与 fixslice}

TWINE 的 4-bit S 盒天然适合 16-byte TBL：一条 TBL 并行完成 16 个 nibble
查找，第二次 TBL 完成轮置换。SM4 的 8-bit S 盒不能直接由单个 16 项表
表示，本项目使用 16 行固定表和掩码选择，访问轨迹与秘密无关；这是常量时间
实现路径，但单块时指令数高。GIFT 的布尔 S 盒和固定 bit permutation
完全不查秘密索引表，更接近 bitslice/fixslice 思路。

对 4-bit S 盒，输入 nibble 本身就是 16 项表的合法索引。ARM NEON
\texttt{TBL} 和 x86 \texttt{PSHUFB} 都能在一个 128-bit 寄存器中同时完成
16 次查找。TWINE 后端先把多个分组的 nibble 排列为固定向量布局，执行 S 盒
后用第二个常量索引向量完成轮置换；下一轮直接消费置换后的布局，避免每轮
保存到内存再重新装载。解密使用逆 S 盒表和逆置换索引，接口语义与参考路径
一致。

SM4 需要处理 8-bit S 盒。实现把输入拆成高、低 nibble，并固定遍历高
nibble 的 16 种可能值：每一行用低 nibble 做 TBL/PSHUFB 查找，再用
“输入高 nibble 是否等于当前行号”的全字节掩码选择结果。循环次数、表行
地址和加载顺序都与秘密无关。四个独立分组经过矩阵转置后同时执行 32 轮，
最后再完成逆转置和大端写回；这样才能摊薄固定扫描 16 行造成的额外指令。

bitslice 则完全取消 S 盒表。GIFT 的四个输入分组先按 bit 位置转置为位平面，
同一条 AND/XOR/NOT 同时作用于多个逻辑 S 盒；固定 bit permutation 通过
掩码交换完成。fixslice 进一步让连续若干轮采用预先选择的数据布局，把一轮
结束后的显式置换吸收到下一轮的 S 盒输入和轮密钥布局中。其代价是装入/写回
时需要转置，因此小消息可能不占优势，而 CTR 等天然存在多个独立计数器的
场景更容易填满并行 lane。

\subsection{新指令集}

\begin{itemize}
  \item \textbf{AES：}ARM AESE/AESD 为本机四路运行路径；x86 AES-NI
  与 VAES 已进入四路/八路实际后端并通过反汇编确认。ARM 的 AESE
  完成 SubBytes、ShiftRows 和 AddRoundKey，随后 AESMC 完成列混合；
  x86 AESENC 把对应步骤合并为一条指令，VAES 则在更宽向量寄存器中同时
  处理多个 128-bit lane。解密分别使用 AESD/AESIMC 与 AESDEC/VAESDEC。
  \item \textbf{GHASH：}PMULL 做两个 64-bit 多项式的 Karatsuba 乘法，
  再以 $x^{128}=x^7+x^2+x+1$ 两次折叠约减；x86 PCLMUL/VPCLMUL
  已进入实际 GHASH 后端。输入在进入指令前转换为 GHASH 的多项式 bit
  次序，三个部分积组合为 256-bit 中间值，约减后再转换回网络字节序。
  四块聚合使四个乘法彼此独立，给 PMULL/CLMUL 留出指令级并行空间。
  \item \textbf{GFNI：}SM4 S 盒按
  $Y=A_1X+C_1,\ S(X)=A_2Y^{-1}+C$ 分解。标量 GFNI 模型穷举 256 个
  输入并逐项等于标准 S 盒，x86 源码用 affine、affine-inverse 和常量异或
  完成多个字节的 S 盒。矩阵和 bit 编号先由标量模型固定，再进入 SIMD
  实现，避免“指令能执行但有限域基底不一致”的隐蔽错误。
  \item \textbf{AES 辅助 SM4：}先用固定扫描预映射到 AES 逆 S 盒值，
  再以 AESE/AESENC 完成四块 SM4 S 盒并逆转 ShiftRows；全程不以秘密值
  直接索引缓存。该路径复用普及度更高的 AES 指令，在没有 SM4 专用指令时
  仍可对四个分组并行处理非线性层，线性层继续由旋转和异或完成。
  \item \textbf{专用 SM4 适配：}ARM SM4E/SM4EKEY 与 Intel
  VSM4RNDS4/VSM4KEY4 均已完成后端适配和静态指令检查。Apple M2 Pro
  未报告 SM4E，Core i9-13900H 未报告 VSM4，因此没有专用 SM4 指令的
  原生性能数据。运行时显式请求该后端会返回
  \texttt{SC\_ERR\_UNSUPPORTED}，而不会落到参考路径后继续标记为硬件。
\end{itemize}

\begin{table}[htbp]
\centering
\caption{优化后端的数据宽度与进入条件}
\begin{tabularx}{\textwidth}{lclX}
\toprule
后端 & 典型批量宽度 & 主要指令/表示 & 进入条件与余数处理\\
\midrule
AES ARM & 4 块 & AESE/AESD、AESMC/AESIMC &
ARM AES feature 可用；不足四块由窄路径完成\\
AES x86 & 4/8 块 & AES-NI、VAES &
CPUID 与 XGETBV 同时允许相应向量状态\\
SM4 shuffle & 4 块 & TBL 或 PSHUFB 固定扫描 &
NEON/SSSE3 可用；尾块回退到较窄实现\\
SM4 GFNI & 4 块 & GF2P8AFFINE(INV) &
x86 GFNI 可用且矩阵模型验证通过\\
GIFT/TWINE & 4 块 & bitslice 或 nibble shuffle &
完整四块走批量内核，余数保持参考语义\\
GHASH & 1/4 块 & PMULL、PCLMUL、VPCLMUL &
有 CLMUL 时执行多项式内核，否则用标量乘法\\
\bottomrule
\end{tabularx}
\end{table}

\section{CTR、GCM 与 XTS}

\subsection{CTR}

CTR 对四类密码全部开放。计数器按分组宽度大端递增，每批生成八个计数器并
调用多块 API；末尾不足一块时只异或所需字节。调用前先验证全部计数器空间，
回绕时不修改输出，并返回错误码 \texttt{SC\_ERR\_COUNTER\_WRAP}。

具体流程是：先在临时副本上递增 \(\lceil len/B\rceil\) 次完成回绕预检；
随后每批复制并递增最多八个计数器，调用分组密码的批量加密生成密钥流，最后
与输入异或。因此加密和解密使用同一函数，也能正确处理原地输入和不足整块
的尾部。

这里的 \(B\) 由分组算法决定：AES/SM4 为 16 字节，GIFT-64/TWINE 为
8 字节；计数器宽度与分组宽度相同，均从最低有效的末字节向前按大端方式
进位。预检使用计数器副本，只有确认整个请求不会回绕后才开始写输出和更新
真实计数器，因此错误返回具有原子性。长度为零时不产生密钥流；尾部只使用
最后一个密钥流块的前若干字节，但真实计数器仍按“已消费一个完整密钥流块”
前进，避免下一次调用重复使用同一 counter block。

批量生成并不会改变 CTR 的安全条件：同一密钥下不得重复使用相同初始计数器。
优化只把连续八个计数器一次交给分组密码，并没有改变计数器取值或异或顺序。
由于密钥流先写入局部数组，输入与输出完全重合时仍会在覆盖输入字节前读取
正确值。

\begin{figure}[htbp]
\centering
\begin{tikzpicture}[node distance=7mm and 8mm]
  \node[flowio] (in) {输入\\counter、in、len};
  \node[flowproc, right=of in] (probe)
    {副本预递增\\$N=\lceil len/B\rceil$ 次};
  \node[flowdecision, right=of probe] (wrap) {回绕？};
  \node[flowerror, above=of wrap] (err)
    {返回回绕错误\\输出和计数器不变};
  \node[flowproc, right=of wrap] (batch)
    {复制并递增\\1--8 个计数器};
  \node[flowproc, below=of batch] (enc)
    {{\scriptsize\texttt{sc\_encrypt\_blocks}}\\生成批量密钥流};
  \node[flowproc, left=of enc] (xor)
    {输入 $\oplus$ 密钥流\\尾批只取所需字节};
  \node[flowdecision, left=of xor] (more) {还有数据？};
  \node[flowio, left=of more] (out)
    {输出结果\\保留递增后 counter};

  \draw[flowarrow] (in) -- (probe);
  \draw[flowarrow] (probe) -- (wrap);
  \draw[flowarrow] (wrap) -- node[right,font=\small]{是} (err);
  \draw[flowarrow] (wrap) -- node[above,font=\small]{否} (batch);
  \draw[flowarrow] (batch) -- (enc);
  \draw[flowarrow] (enc) -- (xor);
  \draw[flowarrow] (xor) -- (more);
  \draw[flowarrow] (more) -- node[above,font=\small]{否} (out);
  \coordinate (loopbottom) at ($(more.south)+(0,-6mm)$);
  \coordinate (loopright) at ($(batch.east)+(8mm,0)$);
  \draw[flowarrow] (more.south) --
    node[left,font=\small]{是} (loopbottom) -| (loopright) -- (batch.east);
\end{tikzpicture}
\caption{CTR 的回绕预检、八路密钥流生成与尾部处理。}
\end{figure}

\begin{lstlisting}[language=C,caption={CTR 的回绕预检和八路批量实现（src/modes.c）}]
blocks = (len + block_size - 1) / block_size;
memcpy(probe, counter, block_size);
for (size_t i = 0; i < blocks; ++i)
    if (!increment_be(probe, block_size))
        return SC_ERR_COUNTER_WRAP;  /* output is untouched */

while (off < len) {
    size_t left = (len - off + block_size - 1) / block_size;
    lanes = left < 8 ? left : 8;
    for (size_t lane = 0; lane < lanes; ++lane) {
        memcpy(counters + lane * block_size, counter, block_size);
        increment_be(counter, block_size);
    }
    sc_encrypt_blocks(ctx, counters, stream, lanes);
    take = len - off < lanes * block_size
         ? len - off : lanes * block_size;
    for (size_t i = 0; i < take; ++i)
        out[off + i] = in[off + i] ^ stream[i];
    off += take;
}
\end{lstlisting}

\subsection{GCM}

GCM 只接受 128-bit 分组的 AES/SM4。96-bit IV 直接形成
$J_0=\mathrm{IV}\|0^{31}\|1$，其他 IV 先经 GHASH。AAD、密文和长度字段
均按 SP~800-38D 处理。解密先以常量时间比较标签；失败返回
\texttt{SC\_ERR\_AUTH} 并清零输出。NIST AES-GCM 与 RFC~8998
SM4-GCM 向量均纳入回归。GHASH 预计算 $H^2,H^3,H^4$，完整的四块组用
四个无数据依赖的乘法聚合。

实现首先计算 \(H=E_K(0^{128})\)。96-bit IV 直接组成
\(J_0=IV\|0^{31}\|1\)，其他长度的 IV 也进入 GHASH。正文数据使用与 CTR
相同的八路批量加密，但只递增计数器低 32 位；认证部分依次吸收 AAD、密文
和二者的 bit 长度，最后计算
\(T=E_K(J_0)\oplus\operatorname{GHASH}_H(A,C)\)。
四块聚合把四次串行更新
\(Y_i=(Y_{i-1}\oplus X_i)H\) 改写为下式：
\[
Y_4=(Y_0\oplus X_1)H^4\oplus X_2H^3\oplus X_3H^2\oplus X_4H .
\]

GHASH 在 \(GF(2^{128})\) 上计算，约减多项式为
\(x^{128}+x^7+x^2+x+1\)。AAD 和密文分别按 16 字节分块，最后不足一块
在右侧补零；末尾再吸收两个大端 64-bit 长度字段，分别记录 AAD bit 数和
密文 bit 数。非 96-bit IV 使用同一 GHASH 规则计算 \(J_0\)，并附加 IV
长度字段。这样快速路径只省略了 IV 的 GHASH，不改变后续 inc32、CTR 或
标签计算。

加密流程先用 \(J_0+1,J_0+2,\ldots\) 的低 32 位计数器产生密文，再对
AAD 和密文认证。解密流程反过来先对收到的密文计算期望标签，并在
\(1\ldots16\) 字节的请求标签长度上做无提前退出的异或归并比较；只有比较
成功才执行 CTR 解密。标签错误、计数器回绕或后端失败都会清零明文输出，
避免调用者误用未认证数据。计数器预检只检查低 32 位，因为 SP~800-38D 的
GCM counter block 使用 inc32，而不是对整个 128-bit \(J_0\) 进位。

四块聚合的关键并非减少乘法次数，而是把原来依赖前一块 \(Y_i\) 的四次乘法
改写为可并行的四项。预计算 \(H^2,H^3,H^4\) 的成本在长消息上被摊薄；
不足 64 字节或四块组后的尾部继续使用单块更新。因此标量、PMULL、
PCLMUL 和 VPCLMUL 后端得到完全相同的 GHASH 值，只在乘法实现和并行宽度
上不同。

\begin{figure}[H]
\centering
\resizebox{0.82\textwidth}{!}{%
\begin{tikzpicture}[node distance=7mm and 12mm]
  \node[flowio] (inputs) {输入\\$K$、IV、AAD、$P/C$};
  \node[flowproc, below left=of inputs] (h)
    {$H=E_K(0^{128})$\\预计算 $H^2,H^3,H^4$};
  \node[flowproc, below right=of inputs] (j0)
    {构造 $J_0$\\96-bit IV 快速路径};
  \node[flowproc, below=of j0] (ctr)
    {低 32 位 inc32\\八路 $E_K(\mathrm{counter})$};
  \node[flowproc, below=of ctr] (cipher)
    {$C=P\oplus$ 密钥流\\加密/解密共用 CTR};
  \node[flowproc, below left=9mm and 2mm of cipher] (ghash)
    {GHASH(AAD, $C$, lengths)\\四块聚合得到 $S$};
  \node[flowproc, below right=9mm and 2mm of cipher] (mask)
    {$E_K(J_0)$\\生成标签掩码};
  \node[flowproc, below=of $(ghash)!0.5!(mask)$] (tag)
    {$T=S\oplus E_K(J_0)$\\支持截断标签};
  \node[flowdecision, below=of tag] (verify) {解密标签相等？};
  \node[flowio, below left=of verify] (ok)
    {是：执行 CTR\\输出明文};
  \node[flowerror, below right=of verify] (bad)
    {否：清零输出\\返回 AUTH};

  \draw[flowarrow] (inputs) -- (h);
  \draw[flowarrow] (inputs) -- (j0);
  \draw[flowarrow] (j0) -- (ctr);
  \draw[flowarrow] (ctr) -- (cipher);
  \draw[flowarrow] (cipher) -- (ghash);
  \draw[flowarrow] (h) |- (ghash);
  \coordinate (maskroute) at ($(mask.east)+(5mm,0)$);
  \draw[flowaux] (j0.east) -| (maskroute) -- (mask.east);
  \draw[flowarrow] (ghash) -- (tag);
  \draw[flowarrow] (mask) -- (tag);
  \draw[flowarrow] (tag) -- (verify);
  \draw[flowarrow] (verify) --
    node[pos=.42,above left,fill=white,inner sep=1pt,font=\small]{是} (ok);
  \draw[flowarrow] (verify) --
    node[pos=.42,above right,fill=white,inner sep=1pt,font=\small]{否} (bad);
\end{tikzpicture}%
}
\caption{GCM 中 CTR 数据路径、四块 GHASH 与先认证后解密流程。}
\end{figure}

\noindent\begin{minipage}{\linewidth}
\begin{lstlisting}[language=C,caption={GCM 的四块 GHASH 聚合（src/modes.c）}]
memcpy(h2, h, 16);  ghash_mul(h2, h);
memcpy(h3, h2, 16); ghash_mul(h3, h);
memcpy(h4, h3, 16); ghash_mul(h4, h);

memcpy(term[0], y, 16);
xor16(term[0], data);         /* Y0 xor X1 */
memcpy(term[1], data + 16, 16);
memcpy(term[2], data + 32, 16);
memcpy(term[3], data + 48, 16);
ghash_mul(term[0], h4);
ghash_mul(term[1], h3);
ghash_mul(term[2], h2);
ghash_mul(term[3], h);
memcpy(y, term[0], 16);
xor16(y, term[1]);
xor16(y, term[2]);
xor16(y, term[3]);
\end{lstlisting}
\end{minipage}

\texttt{ghash\_mul} 在 ARM 上分派到 PMULL，在 x86 上分派到
PCLMUL/VPCLMUL，否则回退到常量时间的逐 bit 标量乘法。解密在产生明文前
先计算并常量时间比较标签；失败时清空输出。

\subsection{XTS}

XTS 使用两个独立密钥上下文，拒绝短于 16 字节的数据单元，并对非整块尾部
执行 ciphertext stealing。IEEE~1619 约定在每个字节内左移并以
\texttt{0x87} 约减；GB/T~17964-2021 约定在每个字节内右移并以
\texttt{0xe1} 约减。OpenSSL~3.6 的 SM4-XTS 默认值为 GB，本项目同时
保留 IEEE 选项。普通分组按八路预生成 tweak 并批量执行 XEX；CTS 在写回前
保存尾部，因而支持原地加解密。

XTS 使用数据密钥 \(K_1\) 和 tweak 密钥 \(K_2\)。先计算
\(T_0=E_{K_2}(i)\)，再对第 \(j\) 个完整分组执行
\(C_j=E_{K_1}(P_j\oplus T_j)\oplus T_j\)，并令
\(T_{j+1}=\alpha T_j\)。实现一次生成八个连续 tweak，把前异或后的八个分组
交给同一个批量后端，然后再异或各自 tweak。

公共函数 \texttt{sc\_xts\_tweak\_from\_data\_unit} 把 64-bit 数据单元号
编码成 16 字节初始输入：IEEE 约定写入低地址八字节并采用小端编码，GB/T
约定写入末八字节并采用大端编码。这个 16 字节值还不是每块直接使用的
tweak，必须先经过 \(E_{K_2}\) 得到 \(T_0\)。初始化检查两个上下文的算法和
分组长度完全一致，并拒绝 GIFT/TWINE 以及未知 tweak 枚举。

普通完整分组采用 XEX 结构：前异或把分组绑定到所在数据单元和块号，数据
密钥完成真正的加/解密，后异或恢复字节表示。八路实现先在局部数组保存八个
tweak，再批量执行前异或、分组密码和后异或；因此解密可以直接选择
\texttt{sc\_decrypt\_blocks}，不需要逐块绕回公共单块接口。IEEE 与 GB/T
的差异只封装在 \texttt{xts\_mul\_x}，其余 XEX 和 CTS 控制流完全共用。

\begin{center}
\centering
\resizebox{0.82\textwidth}{!}{%
\begin{tikzpicture}[node distance=7mm and 8mm]
  \node[flowio] (input) {输入\\$K_1$、$K_2$、单元号 $i$};
  \node[flowproc, right=of input] (t0)
    {$T_0=E_{K_2}(i)$\\加密初始 tweak};
  \node[flowproc, right=of t0] (tweaks)
    {$T_{j+1}=\alpha T_j$\\\texttt{xts\_mul\_x} 八路生成};
  \node[flowproc, below=of tweaks] (pre)
    {前异或\\$P_j\oplus T_j$};
  \node[flowproc, left=of pre] (data)
    {八路数据加/解密\\$E_{K_1}$ 或 $D_{K_1}$};
  \node[flowproc, left=of data] (post)
    {后异或\\结果 $\oplus T_j$};
  \node[flowdecision, below=of post] (tail) {存在非整块尾部？};
  \node[flowio, below left=of tail] (normal)
    {否：输出完整分组};
  \node[flowerror, below right=of tail] (cts)
    {是：保存尾部输入\\执行 CTS 再写回};

  \draw[flowarrow] (input) -- (t0);
  \draw[flowarrow] (t0) -- (tweaks);
  \draw[flowarrow] (tweaks) -- (pre);
  \draw[flowarrow] (pre) -- (data);
  \draw[flowarrow] (data) -- (post);
  \draw[flowarrow] (post) -- (tail);
  \draw[flowarrow] (tail) --
    node[pos=.42,above left,fill=white,inner sep=1pt,font=\small]{否} (normal);
  \draw[flowarrow] (tail) --
    node[pos=.42,above right,fill=white,inner sep=1pt,font=\small]{是} (cts);
\end{tikzpicture}%
}
\captionsetup{hypcap=false}
\captionof{figure}{XTS 的双密钥、八路 tweak/XEX 与 CTS 分支。}
\end{center}

\begin{lstlisting}[language=C,caption={XTS 的八路 tweak/XEX 流程（src/modes.c）}]
uint8_t t[16];
sc_encrypt_block(tweak_ctx, tweak, t);
for (size_t lane = 0; lane < lanes; ++lane) {
    memcpy(tweaks + 16 * lane, t, 16);
    for (size_t i = 0; i < 16; ++i)
        work[16 * lane + i] = input[16 * (block + lane) + i]
                            ^ t[i];
    xts_mul_x(t, convention);
}
sc_encrypt_blocks(data_ctx, work, work, lanes);
for (size_t lane = 0; lane < lanes; ++lane)
    for (size_t i = 0; i < 16; ++i)
        output[16 * (block + lane) + i] =
            work[16 * lane + i] ^ tweaks[16 * lane + i];
\end{lstlisting}

若数据单元含 \(m\) 字节尾部，CTS 先用当前 tweak 得到临时完整密文
\texttt{CC}，把 \texttt{CC} 的前 \(m\) 字节作为尾部密文，再用“尾部明文
\(\|\)\texttt{CC} 剩余字节”和下一个 tweak 加密成倒数第一个完整密文。
代码在覆盖输出前保存尾部输入，因此 17、31、63-byte 等原地加解密不会破坏
尚未读取的数据。

解密 CTS 先用下一个 tweak 解开倒数第一个完整密文，得到由“尾部明文前缀”
和 \texttt{CC} 后缀组成的临时块；再把收到的短尾部密文补到该后缀前方，
使用当前 tweak 解出原倒数第一个完整明文。加密与解密都会在任何写回之前
复制短尾部，解决输入输出相同且写入完整块会覆盖尾部的别名问题。整块长度
时不进入 CTS；恰好 16 字节的数据单元只执行一次普通 XEX。
\section{测试与工程验证}

\begin{table}[htbp]
\centering
\caption{验证层次}
\begin{tabularx}{\textwidth}{lX}
\toprule
层次 & 内容\\
\midrule
KAT & AES 三种密钥、SM4、GIFT-64、TWINE-80/128、NIST GCM、RFC 8998\\
内部交叉 & 参考/表/shuffle/硬件后端，固定种子随机输入，原地与非对齐语义\\
模式边界 & CTR 尾部/回绕、GCM AAD/错误标签清零、XTS CTS/最小数据单元\\
OpenSSL & x86 上 \XOpenSSLChecks 项 AES/SM4 ECB、CTR、GCM、XTS
随机差分；覆盖全部可用后端\\
工具链 & \texttt{-Wall -Wextra -Wpedantic -Werror}、ASan、UBSan\\
ISA & ARM AESE/TBL/PMULL；\newline
x86 AES/VAES/PSHUFB/GFNI/CLMUL；SM4/VSM4 已适配，实验设备未报告支持\\
\bottomrule
\end{tabularx}
\end{table}

\textbf{算法级验证。}每种算法先运行标准 KAT，确保密钥扩展、端序和轮函数
正确；随后用固定种子生成随机密钥与输入，把参考、T-table、shuffle 和所有
可执行 ISA 后端逐块比较。AES 覆盖 128/192/256-bit 三种密钥，SM4 覆盖
三种表布局以及 shuffle、AES 辅助、GFNI，GIFT/TWINE 同时比较单块参考与
四块路径。测试不只验证密文，还用同一后端解密并检查恢复原文。

\begin{lstlisting}[language=C,caption={后端与参考实现的批量交叉验证框架（tests/test\_symcrypto.c）}]
sc_init(&ref, algorithm, SC_BACKEND_REF, key, key_len);
sc_init(&opt, algorithm, requested_backend, key, key_len);

sc_encrypt_blocks(&ref, input, expected, blocks);
sc_encrypt_blocks(&opt, input, actual, blocks);
CHECK(memcmp(actual, expected, blocks * block_size) == 0);

sc_decrypt_blocks(&opt, actual, recovered, blocks);
CHECK(memcmp(recovered, input, blocks * block_size) == 0);
\end{lstlisting}

\textbf{API 与边界验证。}批量测试覆盖 0、1、并行宽度前后以及不能整除 lane
数的分组数量，并分别使用对齐、非对齐和原地缓冲区。CTR 检查四种算法的
零长度、整块、尾部与全 0xff 附近回绕，确认失败时输入、输出和 counter
均不改变。GCM 覆盖空 AAD、任意长度 IV、截断标签、错误标签和失败清零。
XTS 覆盖 16 字节最小单元、17/31/63 字节 CTS、512~B 与 4~KiB 数据单元、
IEEE/GB 两种 tweak 约定以及原地加解密。

\textbf{外部差分与 ISA 验证。}OpenSSL 差分使用独立 API 计算 AES/SM4 的
ECB、CTR、GCM 和 XTS，包含非 96-bit GCM IV 与 XTS-CTS。x86 共完成
\XOpenSSLChecks 项差分检查；运行时专用测试直接进入 AES-NI/VAES、
PSHUFB、GFNI、PCLMUL/VPCLMUL 函数，而不是只测试自动分派。构建阶段再对
实际后端符号反汇编，确认预期指令确实出现在相应函数中。ASan 检查越界和
生命周期错误，UBSan 检查移位、对齐等未定义行为，警告则由
\texttt{-Werror} 提升为构建失败。

\section{实验方法}

ARM 实验平台为 \BenchCPU（\BenchPlatform），编译器为 \BenchCompiler；
x86 实验平台为 \XBenchCPU（\XBenchPlatform），编译器为
\XBenchCompiler，OpenSSL 为 \XOpenSSLVersion。
消息长度取 64~B、512~B、4~KiB、8~KiB、64~KiB、1~MiB。每组预热
3 次、正式运行 15 次；每个样本至少处理 256~KiB，记录完整纳秒值，
汇总 median、Q1、Q3、IQR、ns/byte 和十进制 GB/s。密钥扩展在计时区间外，
模式的计数器、tweak、GHASH 和标签均在计时区间内。volatile checksum
防止编译器删除计算。

每个配置先完成密钥初始化，再执行 3 次不计入统计的预热，使代码页、表和
数据进入稳定缓存状态。正式样本重复处理同一消息长度，内部循环次数按
“总处理量至少 256~KiB”确定，避免 64~B 样本被计时器分辨率支配。每次循环
都会把输出归入 checksum，保证编译器不能把已知输入的加密调用删除。分组
基准分别记录加密和解密；模式基准把 counter/tweak 生成、异或、GHASH、
标签和 CTS 都算入端到端时间。

15 个样本排序后取中位数作为中心值，取 Q1 与 Q3 的差
\(\mathrm{IQR}=Q3-Q1\) 描述短时调度抖动。吞吐率按
\[
  \mathrm{GB/s}=\frac{\text{处理字节数}}{\text{纳秒数}}
\]
计算，使用十进制 GB；\(\mathrm{ns/byte}\) 是其倒数形式，更适合观察低吞吐
算法。报告保留全部样本，不只保存最终中位数：ARM 原始数据位于
\path{results/raw/native_samples.csv}，x86 原始数据位于
\path{results/raw/x86_samples.csv}，对应汇总分别位于
\path{results/summary/native_summary.json} 和
\path{results/summary/x86_summary.json}。

x86-64 结果同时记录单调时钟和序列化 RDTSCP。文中的 TSC ticks/byte
表示 invariant-TSC 计数，不等同于随睿频变化的物理核心周期。ARM 与 x86
结果分别保存，不把一个平台的数据改名为另一个平台。

两台机器的编译器、操作系统、缓存结构和频率策略不同，因此正文只在同一
平台内计算后端相对收益。跨平台图用于确认“某类优化是否同方向有效”，不把
绝对 GB/s 差异直接归因于单条指令。x86 的 CPU 特征、测试计数和 ISA
执行状态另存于 \path{results/summary/x86_status.json}，用于区分“原生执行”
与“仅完成目标 ISA 编译检查”。

\section{结果与分析}

\subsection{ARM64 原生执行结果}

ARM64 的 15 次正式测量样本保存在
\path{results/raw/native_samples.csv}，汇总统计保存在
\path{results/summary/native_summary.json}。本小节的图表和表格均使用
这两份本机结果。

\begin{figure}[htbp]
\centering
\includegraphics[width=\textwidth]{results/figures/block_backends_1m.pdf}
\caption{ARM64 上 1~MiB 分组加密吞吐率，横轴为对数尺度。}
\end{figure}

在 1~MiB 数据上，AES 参考、四表和硬件后端分别为
\AESRefGbps、\AESTableGbps、\AESHWGbps~GB/s。T-table 相对参考为
\AESTableGain$\times$，AESE/AESMC 为 \AESHWGain$\times$。
SM4 四表由 \SMFourRefGbps 提升至 \SMFourTableGbps~GB/s，即
\SMFourTableGain$\times$。1~KiB 单表非常接近四表，2~KiB 重叠读取在
本机反而较慢，说明“更小表”必须与额外旋转/装配成本一起评价。

硬件 AES 加密达到 8.812~GB/s，解密为 8.128~GB/s；两者都由四路
AESE/AESD 内核执行，解密稍低来自逆轮密钥和 AESIMC 数据路径。相比之下，
AES 四表只有 0.199~GB/s，说明在支持专用指令的平台上，查表减少标量逻辑的
收益远小于整轮硬件指令带来的收益。T-table 仍有价值，因为它揭示了没有
密码扩展时“计算换内存”的性能边界。

SM4 的 ARM 结果更能区分不同软件方法：固定访问 shuffle 为
0.112~GB/s，相对 0.066~GB/s 参考路径约提升 1.70 倍；AES 指令辅助路径为
0.096~GB/s，与 4~KiB T-table 的 0.097~GB/s 接近。GIFT 四路 bitslice 从
0.017 提升至 0.069~GB/s，约 4.1 倍；TWINE-80 的 TBL 路径从 0.065 提升
至 0.162~GB/s，约 2.5 倍。这说明轻量密码的主要收益来自多块数据布局，而
不只是把一个标量 S 盒替换为向量查表。

\begin{figure}[htbp]
\centering
\includegraphics[width=\textwidth]{results/figures/size_scaling.pdf}
\caption{不同消息长度下 AES/SM4 后端的稳态吞吐。}
\end{figure}

硬件 AES 在 64~B 时已经领先，随后趋于稳定；本实验把函数调用和循环调度
算入短消息成本。SM4 shuffle 每轮固定扫描 16 行表，并以四块并行摊薄装配
开销；其目标是展示无秘密相关访存，而不是声称单块最高性能。
TWINE TBL 在完整批次上明显快于参考路径，但短消息仍会受到装入、nibble
重排和不足四 lane 的影响，说明向量查表必须与批量数据布局一起评价。

从长度曲线看，64~B 到 512~B 区间仍包含调用、分派和批次不满的固定开销；
达到 4~KiB 后多数曲线趋于稳定。AES 硬件路径较早进入平台期，是因为轮函数
本身很短，随后主要受加载、保存和循环调度限制。SM4 shuffle 需要完整四路
数据才能摊薄转置，短消息上的相对收益因此低于 1~MiB 稳态结果。

\begin{figure}[htbp]
\centering
\includegraphics[width=\textwidth]{results/figures/mode_throughput.pdf}
\caption{512~B 与 4~KiB 的 CTR/GCM/XTS 吞吐率。}
\end{figure}

AES-CTR 能直接发挥硬件块加密能力；XTS 还要执行第二个密钥的 tweak 加密与
每块乘 $\alpha$，因此低于 CTR。GCM 已用 $H^2$--$H^4$ 消除四块组内的
串行依赖，跨组仍保留必要的 GHASH 依赖。IEEE 与 GB 两种 SM4-XTS
吞吐相近，符合两者只改变轻量 tweak 乘法的预期。

在 1~MiB 上，AES-CTR、GCM 加密和 XTS 加密分别为 1.727、0.913 和
1.429~GB/s。CTR 与 8.812~GB/s 原始分组内核之间的差距来自计数器生成、
密钥流保存和逐字节异或；GCM 在此基础上还计算 GHASH，因此约为 CTR 的
一半；XTS 需要两次 tweak 异或和每块乘 \(\alpha\)，但不含 GHASH，所以
位于二者之间。GCM 解密为 0.880~GB/s，与加密接近，表明先认证后解密的
额外读取没有改变主要瓶颈。

\small
\begin{longtable}{lllrrr}
\caption{ARM64 上 1~MiB 全部实测结果}\\
\toprule
算法 & 后端 & 操作 & GB/s & ns/byte & IQR\\
\midrule
\endfirsthead
\toprule
算法 & 后端 & 操作 & GB/s & ns/byte & IQR\\
\midrule
\endhead
\OneMiBResultRows
\bottomrule
\end{longtable}
\normalsize

\subsection{x86-64 原生执行结果}

x86 状态文件记录为 \texttt{passed}。CPUID/XGETBV 检测到的 AES-NI、
SSSE3、AVX2、VAES、GFNI、PCLMUL 与 VPCLMUL 均由专用测试直接调用，
并与参考实现逐项比较。GFNI 硬件路径还穷举全部 256 个 SM4 S 盒输入；
PCLMUL 和 VPCLMUL 则分别对照独立标量 GHASH oracle。Intel VSM4 后端
已经完成适配和静态指令检查，但 Core i9-13900H 的 CPUID 未报告 VSM4
支持，因此显式请求 \texttt{sm4-hw} 返回
\texttt{SC\_ERR\_UNSUPPORTED}，本报告不提供 VSM4 原生吞吐率。

x86 的 15 次正式测量样本保存在
\path{results/raw/x86_samples.csv}，汇总统计保存在
\path{results/summary/x86_summary.json}，CPU 特征、测试状态与检查数量
保存在 \path{results/summary/x86_status.json}。下述图表和表~4 使用这些
x86 实测文件。

x86 的相对趋势与 ARM 一致，但硬件宽度和指令组合不同。AES 四表由
0.137 提升至 0.267~GB/s，约 1.95 倍；AES-NI/VAES 批量后端达到
21.313~GB/s，约为参考实现的 155.6 倍，解密为 16.157~GB/s。这里的
硬件后端数字同时包含真实函数调用、循环和内存读写，不是单条 AESENC 的
理论吞吐。

SM4 上，GFNI 为 0.299~GB/s，是 0.085~GB/s 参考路径的约 3.52 倍；
PSHUFB shuffle 为 0.220~GB/s，约 2.59 倍；AES 辅助路径为
0.174~GB/s，约 2.05 倍。4~KiB、1~KiB、2~KiB T-table 分别为
0.132、0.127、0.093~GB/s，再次显示表越小并不自动越快。GFNI 加密和
解密分别为 0.299 与 0.300~GB/s，符合 SM4 解密只反转轮密钥次序的结构。

GIFT 四路路径由 0.026 提升到 0.103~GB/s，接近 4 倍；TWINE-80
PSHUFB 路径由 0.065 提升到 0.346~GB/s，约 5.32 倍。后者高于 ARM
对应相对收益并不能单独归因于 PSHUFB，因为编译器、调度、缓存和主频均不同；
可以确认的是，两台平台都证明了“向量查表必须配合多块布局”的方向性结论。

\begin{figure}[htbp]
\centering
\includegraphics[width=\textwidth]{results/figures/x86_block_backends_1m.pdf}
\caption{x86-64 上 1~MiB 分组加密吞吐率；数据来自独立 x86 汇总。}
\end{figure}

\begin{figure}[htbp]
\centering
\includegraphics[width=\textwidth]{results/figures/x86_size_scaling.pdf}
\caption{x86-64 上 AES/SM4 后端随消息长度的吞吐变化。}
\end{figure}

\begin{figure}[htbp]
\centering
\includegraphics[width=\textwidth]{results/figures/x86_mode_throughput.pdf}
\caption{x86-64 上 CTR、GCM 与 XTS 的 512~B/4~KiB 吞吐率。}
\end{figure}

x86 的 1~MiB AES-CTR 为 2.129~GB/s，GCM 加/解密为
0.806/0.785~GB/s，XTS 加/解密为 1.230/1.248~GB/s。GCM 的比例下降
说明 GHASH 和认证数据流成为主要开销，不能只凭 21.313~GB/s 的 AES
分组速度推测模式吞吐。SM4 AES 辅助后端在 GCM 中达到约
0.147~GB/s、在 GB-XTS 中达到约 0.159~GB/s，均明显高于参考模式；
模式层批量调用确实把分组后端收益传递到了端到端实现。

\small
\begin{longtable}{lllrrr}
\caption{x86-64 上 1~MiB 全部实测结果}\\
\toprule
算法 & 后端 & 操作 & GB/s & ns/byte & TSC ticks/byte\\
\midrule
\endfirsthead
\toprule
算法 & 后端 & 操作 & GB/s & ns/byte & TSC ticks/byte\\
\midrule
\endhead
\XOneMiBResultRows
\bottomrule
\end{longtable}
\normalsize

\FloatBarrier
\subsection{ARM64 与 x86-64 对照分析}

\textbf{总体趋势。}把两组原生结果放在一起看，可以得到“基础软件路径接近、
x86 专用向量后端优势更大、模式吞吐不与分组吞吐成比例”的总体结论。两组
数据来自不同 CPU、操作系统和编译器，因此下面的比例用于描述本项目的两套
实测实现，而不是严格受控的微架构优劣结论。

\textbf{AES 对比。}AES 标量参考在 ARM64 与 x86-64 上分别为 0.134 和
0.137~GB/s，几乎处于同一水平；四表路径分别为 0.199 和
0.267~GB/s，x86 约高 34\%。进入硬件路径后差距明显扩大：ARM
AESE/AESMC 为 8.812~GB/s，相对本平台参考实现提升 65.72 倍；x86
AES-NI/VAES 为 21.313~GB/s，相对参考实现提升约 155.6 倍，绝对吞吐约为
ARM 的 2.42 倍。

\textbf{SM4 与轻量密码对比。}SM4 的 ARM/x86 参考实现分别为
0.066/0.085~GB/s，4~KiB T-table 为 0.097/0.132~GB/s，shuffle 为
0.112/0.220~GB/s，AES 辅助路径为 0.096/0.174~GB/s。x86 的优势从参考
路径约 1.29 倍扩大到 shuffle 的约 1.96 倍，并且还能使用 ARM 平台没有的
GFNI 路径达到 0.299~GB/s。轻量密码方面，GIFT 批量路径在 ARM/x86 上为
0.069/0.103~GB/s，TWINE-80 为 0.162/0.346~GB/s。两台平台都证明
bitslice 或 shuffle 必须与多块布局结合，而 x86 在 TWINE 的 PSHUFB 路径上
收益更明显。

\textbf{工作模式对比。}模式吞吐没有简单复制原始 AES 的 2.42 倍差距：
AES-CTR 在 ARM/x86 上为 1.727/2.129~GB/s，x86 约高 23\%；
AES-GCM 为 0.913/0.806~GB/s，ARM 反而约高 13\%；AES-XTS 为
1.429/1.230~GB/s，ARM 约高 16\%。这说明 GCM 的 GHASH、XTS 的
tweak/XEX、计数器生成和内存调度会重新决定端到端瓶颈，因此模式性能不能
由分组峰值直接推算。共同结论是：硬件 AES 优势明显；固定访问/bitslice
依赖多块；模式还需优化 GHASH 等外围路径。

\section{安全、性能与平台权衡}

\begin{enumerate}
  \item \textbf{T-table：}通常减少指令，但秘密相关索引会暴露缓存访问模式；
  选择该后端时必须同时评估吞吐率与侧信道风险。
  \item \textbf{Shuffle/bitslice：}访问模式固定，更适合常量时间；需要足够
  多并行块才能摊薄转置和数据装配。
  \item \textbf{硬件密码指令：}AES/PMULL 兼顾性能和固定执行路径，但仍不能
  替代协议级 nonce 管理、密钥生命周期和故障防护。
  \item \textbf{运行时分派：}可移植二进制必须检查 CPUID/XGETBV 或 ARM
  feature；“编译器接受 intrinsic”不等于当前 CPU 可以执行。
  \item \textbf{x86 平台：}AES-NI/VAES、SSSE3、GFNI 与 CLMUL 已完成
  原生执行和性能采样；VSM4 已完成适配，但实验设备不支持该指令集。
\end{enumerate}

\begin{table}[htbp]
\centering
\caption{不同实现路径的性能来源与安全关注点}
\begin{tabularx}{\textwidth}{lXX}
\toprule
路径 & 性能来源 & 需要关注的边界\\
\midrule
标量参考 & 数据表示直接、便于审计与交叉验证 &
指令数多；若直接使用普通 S 盒表仍可能有秘密相关访问\\
T-table & 合并 S 盒、线性层和置换，减少轮内指令 &
索引依赖秘密，缓存/预取轨迹可能泄漏；表尺寸影响 L1\\
Shuffle & 固定表地址，一条指令并行多个 nibble/byte &
需要扫描、转置和足够多 lane；固定访问不等于完整侧信道认证\\
Bitslice/fixslice & 用布尔网络并行多个分组，取消秘密索引 &
装入/写回转置成本高，密钥调度与轮布局更复杂\\
硬件密码指令 & 一条指令完成整轮或有限域乘法，依赖链更短 &
必须正确检测 CPU/OS 状态；仍需保护密钥、nonce 和错误路径\\
\bottomrule
\end{tabularx}
\end{table}

工作模式还引入了与轮函数无关的安全条件。CTR 和 GCM 在同一密钥下不得
重复 counter/IV，否则密钥流复用会直接泄漏明文关系；GCM 标签长度越短，
伪造成功概率越高，调用者还必须限制失败尝试次数。XTS 的 tweak 不是 nonce，
但数据单元号必须稳定且与扇区位置一致；XTS 只提供存储加密的窄块语义，
不提供消息完整性。项目实现回绕检测、标签常量时间比较和失败清零，是为了
保证 API 边界正确，但这些检查不能替代上层协议的密钥轮换和状态管理。

“固定访问”也不等于已经通过侧信道认证。shuffle/bitslice 消除了最明显的
秘密索引表，但实际系统仍可能受到分支预测、编译器重写、共享执行资源、
功耗、电磁和故障注入影响。硬件 AES/GFNI/CLMUL 通常提供更稳定的数据路径，
仍需结合目标处理器文档和部署威胁模型评估。因而本报告把性能、访存模式和
平台可用性分开陈述，不给任何后端附加超出测试证据的安全等级。

\FloatBarrier
\section{结论}

本项目建立了从正确性基线到多类优化和三种工作模式的完整链条。实测说明：
表优化的收益受缓存与额外指令共同约束；硬件 AES 的数量级优势最稳定；
shuffle/fixslice 必须和多块数据布局一起设计；GCM/XTS 不能只优化分组密码
而忽略 GHASH、计数器和 tweak。ARM64 与 x86-64 均已完成原生验证；
x86 结果额外证明 GFNI、VAES、PCLMUL/VPCLMUL 的实际执行，并以独立文件
保存 RDTSCP TSC ticks/byte。VSM4 后端已经完成适配和静态指令检查，但
实验使用的 Core i9-13900H 不支持该指令集，因此没有 VSM4 原生性能数据。

在算法层，AES/SM4 的参考、T-table、shuffle 和硬件路径使用同一密钥语义，
GIFT/TWINE 则证明 bitslice 与 nibble shuffle 可以把轻量密码扩展到多块
并行。T-table 展示了空间换计算的经典方法，但 1/2/4~KiB 方案的实测差异
说明缓存占用、旋转和非对齐装配必须同时考虑。ARM AESE/AESD、x86
AES-NI/VAES、GFNI 以及 PMULL/PCLMUL/VPCLMUL 的原生结果进一步说明，
专用指令的收益只有在批量内核、正确轮密钥布局和真实运行时分派共同成立时
才能转化为端到端吞吐。

在模式层，CTR 通过八路计数器直接复用多块加密；GCM 把八路 CTR 与四块
GHASH 聚合组合，并保持任意 IV、AAD、截断标签和先认证后解密语义；XTS
把八路 tweak/XEX、IEEE/GB 两种乘法约定和原地安全 CTS 统一到双上下文
接口。结果证明，最快的分组密码后端并不自动得到同等比例的 GCM/XTS 提升，
因为 GHASH、异或、计数器、tweak 和尾部控制流会成为新的瓶颈。

最终实现同时保留 ARM64 与 x86-64 的原始样本、汇总统计、平台特征和
指令检查证据。可执行的后端已经完成 KAT、随机交叉、OpenSSL 差分、
sanitizer 和原生性能测试；已适配但设备不支持的 SM4E/VSM4 则明确停留在
编译与静态指令检查层级。由此，报告中的每项功能、优化方法和性能结论都能
对应到具体代码路径与相应验证证据。

\clearpage
\begin{thebibliography}{9}\sloppy
\bibitem{fips197} NIST, \emph{FIPS 197: Advanced Encryption Standard}, 2023.
\url{https://csrc.nist.gov/pubs/fips/197/final}
\bibitem{sm4} 国家标准化管理委员会，\emph{GB/T 32907-2016 SM4 分组密码算法}.
\url{https://openstd.samr.gov.cn/bzgk/std/newGbInfo?hcno=7803DE42D3BC5E80B0C3E5D8E873D56A}
\bibitem{gift} Banik et al., \emph{GIFT: A Small Present}, CHES 2017.
\url{https://www.iacr.org/archive/ches2017/10529227/10529227.pdf}
\bibitem{twine} Suzaki et al., \emph{TWINE: A Lightweight Block Cipher for Multiple Platforms}, SAC 2012.
\url{https://jpn.nec.com/rd/tg/dss/pdf/twine_SAC_full_v5.pdf}
\bibitem{modes} NIST, \emph{Block Cipher Modes}.
\url{https://csrc.nist.gov/projects/block-cipher-techniques/bcm}
\bibitem{rfc8998} P. Yang, \emph{RFC 8998: ShangMi Cipher Suites for TLS 1.3}, 2021.
\url{https://www.rfc-editor.org/rfc/rfc8998.html}
\bibitem{gfni} W. Guo, \emph{Efficient Constant-Time Implementation of SM4}, 2022/1154.
\url{https://eprint.iacr.org/2022/1154.pdf}
\bibitem{opensslxts} OpenSSL, \emph{EVP\_EncryptInit: xts\_standard}.
\url{https://docs.openssl.org/3.4/man3/EVP_EncryptInit/}
\end{thebibliography}

\end{document}
